\section{口语理解}
\label{sec:nlp-speech-understanding}

用户说“后天三个人”，系统要得到的是“修改预约，日期某天，人数3”这样的语义，而不一定是完整转写。\term{口语理解}（spoken language understanding, SLU）从声音及允许的对话历史恢复意图、参数与交流作用。第\ref{sec:nlp-dialogue-intent-crf}小节已在文字上解决了意图分类与槽位标注；换成声音输入后，又多出四项挑战。其一是识别错误传播：识别器把“后天”听成“后台”，后面的文本模型再准确也找不到日期。其二是口语现象：更正、重复和省略在书面语中少见，在口语里却很常见。其三是语境依赖：同一句“可以”在不同历史和语气下作用不同。其四是监督稀缺：同时带声音和语义标注的数据远少于纯文本标注。

方法谱系按“语义模型看到什么输入”推进。最简单的流水线把识别一最佳结果交给第\ref{chap:nlp-question-answering-dialogue}章的文本理解模型；为缓解错误传播，第二类方法把识别的多个候选或词格交给理解模型，让语义证据参与选择文字；第三类端到端模型直接读取声学编码，用语义目标训练，跳过文字接口；最后再把对话历史和韵律加入输入，判断交流作用。SLURP把声音、转写与场景、动作、实体标注配对，可用于分别检查识别误差与语义误差；它为评价组织提供依据，并不是下面各算法的出处\scite{nlp-speech-bastianelli2020}。全节的评价都把单个字段正确与完整记录正确分开：槽位F1高不等于每条请求都被正确执行。

\begin{table*}[htbp]
  \centering
  \begin{tabularx}{\linewidth}{@{}>{\raggedright\arraybackslash}p{28mm}>{\raggedright\arraybackslash}X>{\raggedright\arraybackslash}X>{\raggedright\arraybackslash}X@{}}
    \toprule
    方法 & 解决的问题 & 核心机制 & 新代价\\
    \midrule
    一最佳流水线 & 复用文本意图与槽位模型 & 识别一最佳$\to$联合意图—CRF模型 & 识别错误直接传入语义\\
    候选重排与词格输入 & 一最佳丢失正确文字 & 识别分数与语义分数联合选择；理解模型读取带后验的词格 & 需要识别器导出候选；计算随候选数增长\\
    直接声音分类 & 语义只需少量类别时跳过转写 & 自监督编码器$\to$池化$\to$softmax & 不给出参数；需声音—意图配对数据\\
    语义记录生成 & 参数、更正与撤回需要结构化输出 & 声学编码器$\to$交叉注意力解码器$\to$受约束记录 & 配对标注稀缺；证据定位需另设监督\\
    历史与韵律条件分类 & 同一文字的交流作用不同 & 拼接历史编码；韵律统计或声学表示与文本融合 & 标注一致性低；需按说话者和会话留出\\
    \bottomrule
  \end{tabularx}
  \caption{口语理解方法对照。自上而下，语义模型看到的输入从单一文字扩展到候选集合、声学序列，再到历史与韵律；每一步解决前一步丢失的信息，也带来新的数据或计算需求。}
  \label{tab:nlp-speech-slu-methods}
\end{table*}
\FloatBarrier

\subsection{意图、槽位与口语事实}
\label{sec:nlp-speech-semantics}

本小节另设一个与读书会无关的预约请求：2025年4月17日，用户说“明天两个人，改成后天三个人”。当前有效值是日期4月19日、人数3；4月18日与人数2是被撤回的旧值，需要另行记录，不能和有效值并列。意图说明请求做什么，槽位记录参数；每个槽位再分别保存规范值和证据区间，前者供业务使用，后者供核查。下面先介绍转写后的联合模型及其口语侧补救，再介绍直接从声音预测语义的模型。

\subsubsection{意图分类与CRF槽位标注}
\label{sec:nlp-speech-intent-crf}

文本侧的模型可以直接沿用。第\ref{sec:nlp-dialogue-jointbert}小节的联合模型让一个预训练文本编码器同时服务两个头：句级表示接softmax预测意图$c$，逐位置表示给BIO标签打分；在槽位头上加线性链CRF，可学习标签转移并排除句首\code{I-DATE}一类非法路径。联合目标为
\[
  L=-\log p(c\mid y)-\lambda\log p_{\mathrm{crf}}(z\mid y),
\]
其中$y$是输入文字，$z$是BIO标签序列，$\lambda$由开发集选定。CRF的定义、梯度与Viterbi解码见第\ref{sec:nlp-classical-crf}小节，合法转移掩码见第\ref{sec:nlp-analysis-paradigm-tagging}小节。对“明天／两／个／人／改成／后天／三／个／人”，模型标出两个日期跨度和两个人数跨度；“改成”触发更正关系，把后一组设为有效值、前一组标为撤回。日期按咨询时点解析，“后天”才得到4月19日。

口语侧的新问题在于$y$不再是人工写下的文字，而是识别结果。设识别器输出“明天两个人改成后台三个人”，CRF完全合法地给“后台”标O，日期槽位就缺失了。理解模型在人工转写上训练，从未见过这类替换，错误被原样传播。要定位误差来源，可以把同一语义模型分别接三种输入：人工转写给出语义模型自身的上限；识别一最佳给出流水线的实际表现；两者之差就是识别误差造成的损失。

第三种输入是识别候选集合，让语义证据参与选择文字。设识别器输出$N$个候选$y^{(1)},\ldots,y^{(N)}$，识别对数概率为$\log p_{\mathrm{asr}}(y^{(k)}\mid\bm{X})$，其中$\bm{X}$是声学特征序列。联合重排在候选与语义结果上同时取最大：
\begin{equation}
  (\hat k,\hat c,\hat z)=\argmax_{k,\,c,\,z}\;
  \mu\log p_{\mathrm{asr}}(y^{(k)}\mid\bm{X})
  +\log p(c,z\mid y^{(k)}),
  \label{eq:nlp-speech-slu-rerank}
\end{equation}
$\mu$平衡两种分数。教学示例取三条候选，识别对数概率为$-1.0$（“……后台三个人”）、$-1.4$（“……后天三个人”）、$-2.0$（“……后天四个人”）；语义模型对各自最优记录的对数概率为$-3.0$（日期槽缺失，记录不完整）、$-0.5$、$-0.6$。取$\mu=1$，三者总分为$-4.0$、$-1.9$、$-2.6$，重排选中第二条，而一最佳是第一条。这组数只用于说明计算；候选里有正确文字也不等于一定被选中，评价时要分别统计“候选含正确答案”与“最终选中正确答案”的比例。

候选列表仍有冗余：$N$条候选往往只在一两个词上不同。识别器解码时内部保留的\term{识别词格}（ASR lattice）把这些候选压缩成有向无环图，弧上是词及其后验概率；与第\ref{sec:nlp-analysis-pretrained-lattice}小节的中文分词词格结构相同，只是弧的来源是声学解码而不是词典。LatticeRNN把循环单元推广到词格上，每个节点汇合所有入弧的状态，并按弧后验加权\scite{nlp-speech-ladhak2016}；Transformer编码器也可以用可达关系构造注意力掩码，让每条弧只关注与它同在某条路径上的弧。词格输入保留了全部候选，计算量只随弧数增长；代价是需要识别器导出词格，且训练数据要用同一识别器解码得到，换识别器后分布会变化。另一种更简单的补救是直接在识别输出上训练语义模型，让它见过真实的识别错误。

无论接哪种输入，验收都不能止于槽位F1。标签序列完全合法时，“后天”仍可能被规范为4月18日，或“两个人”没有被撤回。应分别核查意图、跨度、规范值、更正关系和完整记录；SLURP为此还提出了容忍轻微转写差异的实体评分，用于把识别拼写误差与语义误差分开\cite{nlp-speech-bastianelli2020}。完整记录正确率要求一条请求的所有字段同时正确，与业务结果最接近。

\subsubsection{直接声音分类与语义记录生成}
\label{sec:nlp-speech-direct-slu}

流水线中识别器按逐字错误训练，把“后天”错成“后台”和把“的”错成“得”付出同样代价，而对语义而言前者是致命错误。端到端口语理解改用语义目标直接训练读取声音的模型，让对语义重要的声学差异得到更多关注。Serdyuk等最早演示了不经文字中间表示、直接从声音预测意图的模型\scite{nlp-speech-serdyuk2018}；在自监督语音表示出现后，这类模型通常以预训练编码器为主干。

先看只需判断意图的情形。设录音4秒，经第\ref{sec:nlp-speech-ssl}小节的自监督编码器（如wav2vec 2.0或HuBERT）得到每20毫秒一帧的表示，共$T=200$帧。按该小节末的使用方式，对编码器各层输出学一组softmax权重加权求和，得到$\bm{H}\in\mathbb{R}^{T\times d}$，$d=768$。意图是整句的属性，还要把$T$帧汇成一个向量，常用注意力池化：
\begin{equation}
  \bm{h}_{\mathrm{audio}}=\sum_{t=1}^{T}a_t\bm{h}_t,
  \qquad
  a_t=\frac{m_t\exp(\bm{q}^{\top}\bm{h}_t)}{\sum_{t'}m_{t'}\exp(\bm{q}^{\top}\bm{h}_{t'})}.
  \label{eq:nlp-speech-slu-pooling}
\end{equation}
$\bm{h}_t$是$\bm{H}$的第$t$行；$m_t\in\{0,1\}$是有效帧掩码，排除批内补齐的帧；$\bm{q}\in\mathbb{R}^{d}$是可学习查询向量。所有$a_t$相等时注意力池化退化为掩码均值池化。得到的$\bm{h}_{\mathrm{audio}}\in\mathbb{R}^{768}$经线性层和softmax输出$|\mathcal{C}|$个意图的概率，以交叉熵训练。数据较少时冻结编码器、只训练层权重、查询向量与分类头；数据较多时整体微调。

日期和人数还需要逐字段输出，单个池化向量装不下这些信息。语义记录生成保留整条声学序列$\bm{H}\in\mathbb{R}^{T\times d}$作为编码，接一个Transformer解码器，经式\eqref{eq:nlp-llm-cross-attention}的交叉注意力读取声音，逐词元生成序列化的语义记录$r$，以$-\sum_j\log p(r_j\mid\bm{H},r_{<j})$训练。上述请求的目标记录可以是
\begin{quote}
\code{intent=修改预约; date=2025-04-19; party=3;}\\
\code{retract=(date:2025-04-18, party:2)}
\end{quote}
解码时用第\ref{sec:nlp-inference-grammar}小节的文法约束只允许合法字段名、类型和结束位置，生成后再解析检查必填字段、日期范围与更正状态。一种常用变体让解码器先输出转写、再输出记录，并在编码器上加CTC辅助损失；这样仍保留文字用于核查，但转写与语义在同一模型中联合优化，语义部分可以利用声学证据纠正转写，而不是只能接受一最佳。

端到端模型的瓶颈是数据。声音—语义配对标注昂贵，常见的补救有两条：用预训练的自监督或识别编码器提供声学知识，只用少量配对数据学习语义头；用语音合成把纯文本语义数据转成声音，扩充训练集。合成声音的音色和噪声分布与真实录音不同，需与真实来源分开记录、分开报告。

与流水线相比，端到端模型避开了文字接口上的硬决策，延迟也只来自一个模型；代价是中间结果不可读，出错时难以判断是没听清还是没理解。若要求输出证据区间，必须说明监督对应声学帧区间还是转写字符区间，没有转写时不能称为文字跨度。比较两条路线时应固定声音、语义模式和训练资源，报告关键槽位与完整请求正确率。

\subsection{话语行为与情绪表达}
\label{sec:nlp-speech-pragmatics}

“可以”可能表示确认、许可或勉强接受。交流作用由话语行为刻画，其文本侧建模（HMM与CRF序列标注、分层对话编码）已在第\ref{sec:nlp-analysis-dialogue-acts}小节展开；情绪标签则按协议描述表达出的类别或唤醒度、效价量表，不等于直接测量心理状态。声音带来的新问题是：可用证据多了语调与节奏，文字又可能来自有误的识别。实验需要先固定输入范围（声音、文字、历史、说话者、回应对象各取哪些），标注协议允许不确定或多标签并报告标注者一致性，划分时按说话者和会话留出，以免模型记住某人的声音。类别不均衡时报告宏平均F1，量表任务另报告误差与相关性。

\subsubsection{历史条件下的话语行为分类}
\label{sec:nlp-speech-dialogue-act}

一个可复现的输入是“前轮内容、当前文字、当前说话者、回应对象、已确认状态”，用拼接历史的预训练编码器或第\ref{sec:nlp-analysis-dialogue-acts}小节的分层编码器（在线使用时改为单向）做分类，以交叉熵学习请求、告知、确认、否定、更正等标签；一句话兼有确认与补充告知时，改为多个二元标签并在开发集上定阈值。前轮是“周五下午两点方便吗”时，“可以”标为确认时间；前轮是“我可以修改人数吗”时，同样的回答表示许可。只看当前句的分类器面对完全相同的输入，无法区分两者，加入前轮与回应对象后依据才充分。

口语场景下历史本身也来自识别，前轮若被识别错误或被截断，分类器应输出不确定，系统据此发起澄清，例如“你是确认时间，还是允许修改人数”。评价把标签正确与对话后果分开：行为判为“确认”仍可能确认了错误的值，状态更新还需核对被确认的是哪一项。

\subsubsection{韵律特征、SVM与声文融合}
\label{sec:nlp-speech-prosody-classification}

同一句“可以”读得平缓或急促，节奏可能帮助解释表达。\term{韵律}指音高变化、能量、时长与停顿等跨越单个音素的组织。可复现的基线先从有效语音帧提取基频$F_0$、对数能量与倒谱特征，再统计均值、标准差、分位数、斜率、总时长与停顿比例，形成固定长度向量$\bm{v}$。无声段没有基频，应标为缺失并另记有声帧比例，不能按0\,Hz参与平均，否则停顿多的句子会被误算成音调低。

以二分类标签$z_i\in\{-1,1\}$为例，标准化参数只在训练集估计，支持向量机（support vector machine, SVM）最小化
\[
  \tfrac12\lVert\bm{w}\rVert_2^2+C\sum_i\max\bigl(0,\,1-z_i(\bm{w}^{\top}\bm{v}_i+b)\bigr),
\]
$C$由开发集选定\scite{cortes1995}。手工特征维数低、可解释，适合作为基线。更强的做法用式\eqref{eq:nlp-speech-slu-pooling}把自监督声学表示池化为$\bm{h}_{\mathrm{audio}}$，再与文本表示$\bm{h}_{\mathrm{text}}$融合。早融合把两者拼接为$[\bm{h}_{\mathrm{audio}};\bm{h}_{\mathrm{text}}]$后接分类头，可以学习声文交互；晚融合分别训练两个分类器，再以$\alpha\log p_{\mathrm{audio}}+(1-\alpha)\log p_{\mathrm{text}}$组合，某一路缺失时仍能工作；还可以让文字位置经交叉注意力读取声音序列。

教学对照固定“可以”的文字，只把时长从0.8\,s改为0.4\,s并改变音高走势，观察分类器输入怎样变化，而不预设短促就表示不耐烦。只有在留出的说话者与会话上比较纯文字、纯声音与融合模型，才能判断声音提供了多少增益；识别出某种情绪表达后，如何回应仍由对话目标决定。
